\section{Экономическое обоснование разработки и реализации программного средства}

Четвертый раздел посвящен экономическому анализу разработки программного средства сквозной аналитики для \eng{e-commerce} проектов. Хотя дипломный проект является учебным и не преследует коммерческих целей, проведение экономической оценки демонстрирует практическую целесообразность программного решения для гипотетического внедрения. Анализ включает расчет инвестиций в разработку, условную оценку эффекта от автоматизации и показатели экономической эффективности внедрения.

Важно подчеркнуть, что все цифры и расчеты условны и предназначены для демонстрации экономического обоснования в рамках учебного проекта. Система представляется не как коммерческий стартап с реальными продажами, а как инструмент автоматизации, который мог бы сократить временные затраты специалиста и снизить риск ошибок при работе с аналитикой e-commerce.

\subsection{Характеристика программного средства}

Разработанное программное средство предназначено для автоматизации процессов управления сквозной аналитикой в отделах digital-маркетинга и \eng{SEO}-агентств. Целевая аудитория включает \eng{SEO}-специалистов, менеджеров проектов, \eng{CRM}-администраторов и владельцев интернет-магазинов, которые нуждаются в централизованной системе для отслеживания источников трафика, рассчета ключевых показателей и формирования аналитических отчетов.

Основные возможности системы:

\begin{enumerate}
    \item Регистрация и управление заявками клиентов на подключение сквозной аналитики.
    \item Создание и ведение проектов \eng{e-commerce} с указанием сайта, ниши, регионального охвата и других параметров.
    \item Загрузка и валидация данных о продажах, событиях, возвратах и источниках трафика в табличном виде.
    \item Диагностика качества подключенной аналитики и идентификация проблем в отслеживании.
    \item Настройка модели аналитики: сопоставление источников трафика, целевых событий и этапов воронки.
    \item Автоматический расчет ключевых показателей: переходы, заказы, выручка, конверсия, средний чек, возвраты, эффективность источников.
    \item Формирование аналитических отчетов и рекомендаций по оптимизации источников трафика.
    \item Управление пользователями, ролями, справочниками и тестовыми данными для демонстрации.
\end{enumerate}

Система может быть внедрена в качестве самостоятельного инструмента внутри агентства или организации либо предложена клиентам через модель \eng{SaaS} (Software as a Service) с оплатой по подписке. Для целей экономического анализа рассматривается вариант внедрения в отдельное подразделение, где система ускоряет выполнение аналитических работ и снижает нагрузку на специалиста.

\subsection{Расчет инвестиций в разработку программного средства}

Инвестиции в разработку включают трудозатраты специалистов, социальные отчисления, инфраструктурные расходы и прочие издержки. Основная составляющая — заработная плата разработчиков и аналитиков, задействованных в проекте.

Для реализации дипломного проекта была сформирована учебная команда, состоящая из одного \eng{Full Stack}-разработчика, одного специалиста по дизайну и \eng{UX}, а также одного технического писателя для документации. В условиях реального коммерческого внедрения такой стек может быть расширен, но для учебных целей он достаточен.

Расчет затрат на основную заработную плату команды разработчиков приведен в таблице~\ref{tab:salary-calculation}.

{\small
\setlength{\tabcolsep}{3pt}
\begin{longtable}{|P{0.30\textwidth}|P{0.15\textwidth}|P{0.15\textwidth}|P{0.15\textwidth}|P{0.15\textwidth}|}
    \caption{Расчет затрат на основную заработную плату команды разработчиков}
    \label{tab:salary-calculation}\\
    \hline
    Специалист & Месячная зарплата, руб. & Кол-во месяцев & Трудозатраты, чел.-мес. & Сумма, руб. \\
    \hline
    \endfirsthead
    \caption*{Продолжение таблицы \thetable}\\
    \hline
    Специалист & Месячная зарплата, руб. & Кол-во месяцев & Трудозатраты, чел.-мес. & Сумма, руб. \\
    \hline
    \endhead
    \eng{Full Stack}-разработчик (backend \eng{Go}, frontend \eng{React}) & 120~000 & 5 & 5 & 600~000 \\
    \hline
    Специалист по дизайну и \eng{UX} & 80~000 & 2 & 2 & 160~000 \\
    \hline
    Технический писатель и аналитик требований & 70~000 & 1 & 1 & 70~000 \\
    \hline
    \textbf{Итого:} & & & \textbf{8} & \textbf{830~000} \\
    \hline
\end{longtable}
}

Социальные отчисления (взносы в пенсионный фонд, фонд социального страхования и другие) составляют в среднем 30\% от фонда оплаты труда. Прочие расходы включают облачные сервисы для размещения тестового сервера, подписки на инструменты разработки, лицензии на \eng{IDE} и другие издержки.

Смета затрат на разработку программного средства приведена в таблице~\ref{tab:development-expenses}.

{\small
\setlength{\tabcolsep}{3pt}
\begin{longtable}{|P{0.50\textwidth}|P{0.20\textwidth}|P{0.20\textwidth}|}
    \caption{Смета затрат на разработку программного средства}
    \label{tab:development-expenses}\\
    \hline
    Статья расходов & Объем & Сумма, руб. \\
    \hline
    \endfirsthead
    \caption*{Продолжение таблицы \thetable}\\
    \hline
    Статья расходов & Объем & Сумма, руб. \\
    \hline
    \endhead
    Основная заработная плата разработчиков и специалистов & 8 чел.-мес. & 830~000 \\
    \hline
    Социальные отчисления (30\% от ФОТ) & — & 249~000 \\
    \hline
    Облачные серверы и инфраструктура для разработки & 5 месяцев & 50~000 \\
    \hline
    Лицензии на инструменты разработки и тестирования & Период разработки & 25~000 \\
    \hline
    Прочие оборотные расходы (материалы, связь, командировки) & — & 21~000 \\
    \hline
    \textbf{Итоговые инвестиции:} & & \textbf{1~175~000} \\
    \hline
\end{longtable}
}

Таким образом, общая сумма инвестиций в разработку программного средства составляет \textbf{1~175~000 рублей}. Эта цифра включает все основные статьи расходов на разработку, включая зарплаты, отчисления и инфраструктуру.

\subsection{Расчет экономического эффекта от реализации программного средства}

Экономический эффект от внедрения системы определяется сокращением времени, необходимого \eng{SEO}-специалисту для выполнения аналитических работ, а также снижением вероятности ошибок при ручной обработке данных.

При ручном выполнении аналитических работ специалист тратит в среднем следующее время на одного клиента в месяц:

\begin{itemize}
    \item Загрузка и очистка данных о продажах, событиях и возвратах из разных источников: \textbf{8 часов}.
    \item Проверка качества и валидация данных для выявления пропусков и аномалий: \textbf{6 часов}.
    \item Ручной расчет показателей конверсии, выручки по источникам, \eng{ROI} и других \eng{KPI}: \textbf{10 часов}.
    \item Подготовка отчета, анализ тенденций и формирование рекомендаций: \textbf{8 часов}.
    \item Прочие операции (встречи, согласование, резервные копии): \textbf{3 часа}.
\end{itemize}

\textbf{Итого в месяц на одного клиента (ручной процесс):} \textbf{35 часов}.

При использовании разработанной системы время специалиста сокращается за счет автоматизации:

\begin{itemize}
    \item Загрузка данных осуществляется через табличный интерфейс с валидацией: \textbf{2 часа} (вместо 8).
    \item Проверка качества данных автоматизирована, специалист только проверяет рекомендации системы: \textbf{1.5 часа} (вместо 6).
    \item Расчет показателей выполняется автоматически, специалист анализирует результаты: \textbf{3 часа} (вместо 10).
    \item Формирование отчета автоматизировано, специалист добавляет рекомендации на основе результатов: \textbf{4 часа} (вместо 8).
    \item Прочие операции: \textbf{1 час} (вместо 3).
\end{itemize}

\textbf{Итого в месяц на одного клиента (с системой):} \textbf{11.5 часов}.

\textbf{Сокращение времени на одного клиента:} $35 - 11.5 = 23.5$ \textbf{часов в месяц}.

Предположим, что специалист обслуживает в среднем \textbf{10 активных проектов} в месяц. При 8-часовом рабочем дне и 22 рабочих днях в месяц (176 часов), сокращение времени приводит к:

\textbf{Общее сокращение времени:} $23.5 \times 10 = 235$ \textbf{часов в месяц} или примерно $235 / 176 \approx 1.3$ \textbf{специалиста}.

Денежный эквивалент этой экономии при стоимости одного специалиста в 120~000 рублей в месяц составляет:

$$\text{Эффект в месяц} = 120~000 \times 1.3 = 156~000 \text{ рублей}.$$

$$\text{Годовой эффект} = 156~000 \times 12 = 1~872~000 \text{ рублей}.$$

Кроме того, внедрение системы снижает вероятность ошибок при расчетах на \textbf{90\%}. Ошибки в аналитике могут привести к некорректным рекомендациям и потере доверия клиента. Условно, стоимость устранения одной критической ошибки и восстановления отношений с клиентом составляет \textbf{20~000 рублей}. Если в году без системы происходит в среднем \textbf{5 критических ошибок} на 10 проектах, то:

$$\text{Эффект от снижения ошибок} = 20~000 \times 5 \times 0.9 = 90~000 \text{ рублей в год}.$$

\textbf{Общий годовой эффект:} $1~872~000 + 90~000 = 1~962~000$ \textbf{рублей}.

\subsection{Расчет показателей экономической эффективности разработки и реализации программного средства}

На основе инвестиций и эффекта рассчитаны показатели экономической эффективности внедрения системы.

\subsubsection*{1. Период окупаемости}

$$\text{Период окупаемости} = \frac{\text{Итоговые инвестиции}}{\text{Годовой эффект}} = \frac{1~175~000}{1~962~000} \approx 0.6 \text{ года} \approx \textbf{7 месяцев}.$$

Это означает, что инвестиции в разработку окупятся за счет экономии времени и снижения ошибок в течение приблизительно 7 месяцев внедрения.

\subsubsection*{2. Индекс прибыльности (Profitability Index)}

$$\text{Индекс прибыльности} = \frac{\text{Годовой эффект}}{\text{Итоговые инвестиции}} = \frac{1~962~000}{1~175~000} \approx 1.67.$$

Значение выше 1 свидетельствует о целесообразности инвестиций. На каждый рубль инвестиций система приносит 1.67 рублей эффекта в первый год.

\subsubsection*{3. Чистая текущая стоимость за три года}

При условии, что в последующие годы эффект остается на уровне первого года, чистая текущая стоимость за три года (без дисконтирования) составляет:

$$\text{NPV}_{\text{3 года}} = (1~962~000 \times 3) - 1~175~000 = 5~886~000 - 1~175~000 = 4~711~000 \text{ рублей}.$$

При учете дисконтирования на 10\% в год значение будет несколько меньше, но остается значительным.

\subsubsection*{4. Рентабельность инвестиций (Return on Investment, ROI)}

$$\text{ROI} = \frac{\text{Чистая прибыль}}{\text{Инвестиции}} \times 100\% = \frac{1~962~000 - 1~175~000}{1~175~000} \times 100\% = \frac{787~000}{1~175~000} \times 100\% \approx 67\%.$$

Положительное значение \eng{ROI} подтверждает финансовую целесообразность проекта.

\subsubsection*{Выводы по экономическому обоснованию}

Проведенный анализ демонстрирует, что программное средство сквозной аналитики для \eng{e-commerce} проектов, несмотря на условный характер расчетов, показывает положительные экономические показатели:

\begin{enumerate}
    \item Инвестиции окупаются за 7 месяцев благодаря сокращению времени специалиста и снижению ошибок.
    \item Индекс прибыльности (1.67) свидетельствует о привлекательности проекта для внедрения.
    \item За три года система может принести чистый результат в 4.7 млн рублей.
    \item \eng{ROI} в 67\% за первый год указывает на эффективность инвестиций.
\end{enumerate}

Таким образом, система представляет практическую ценность не только как учебный проект, но и как потенциальный инструмент для реального внедрения в агентствах, отделах аналитики и цифрового маркетинга. Дальнейшее развитие включает интеграцию с внешними \eng{API} \eng{e-commerce}-платформ, расширение функций аналитики и адаптацию под специфику конкретных рынков и регионов.
